% Options for packages loaded elsewhere
\PassOptionsToPackage{unicode}{hyperref}
\PassOptionsToPackage{hyphens}{url}
\documentclass[
  ignorenonframetext,
  aspectratio=169,
]{beamer}
\newif\ifbibliography
\usepackage{pgfpages}
\setbeamertemplate{caption}[numbered]
\setbeamertemplate{caption label separator}{: }
\setbeamercolor{caption name}{fg=normal text.fg}
\beamertemplatenavigationsymbolsempty
% remove section numbering
\setbeamertemplate{part page}{
  \centering
  \begin{beamercolorbox}[sep=16pt,center]{part title}
    \usebeamerfont{part title}\insertpart\par
  \end{beamercolorbox}
}
\setbeamertemplate{section page}{
  \centering
  \begin{beamercolorbox}[sep=12pt,center]{section title}
    \usebeamerfont{section title}\insertsection\par
  \end{beamercolorbox}
}
\setbeamertemplate{subsection page}{
  \centering
  \begin{beamercolorbox}[sep=8pt,center]{subsection title}
    \usebeamerfont{subsection title}\insertsubsection\par
  \end{beamercolorbox}
}
% Prevent slide breaks in the middle of a paragraph
\widowpenalties 1 10000
\raggedbottom
\AtBeginPart{
  \frame{\partpage}
}
\AtBeginSection{
  \ifbibliography
  \else
    \frame{\sectionpage}
  \fi
}
\AtBeginSubsection{
  \frame{\subsectionpage}
}
\usepackage{iftex}
\ifPDFTeX
  \usepackage[T1]{fontenc}
  \usepackage[utf8]{inputenc}
  \usepackage{textcomp} % provide euro and other symbols
\else % if luatex or xetex
  \usepackage{unicode-math} % this also loads fontspec
  \defaultfontfeatures{Scale=MatchLowercase}
  \defaultfontfeatures[\rmfamily]{Ligatures=TeX,Scale=1}
\fi
\usepackage{lmodern}
\ifPDFTeX\else
  % xetex/luatex font selection
\fi
% Use upquote if available, for straight quotes in verbatim environments
\IfFileExists{upquote.sty}{\usepackage{upquote}}{}
\IfFileExists{microtype.sty}{% use microtype if available
  \usepackage[]{microtype}
  \UseMicrotypeSet[protrusion]{basicmath} % disable protrusion for tt fonts
}{}
\makeatletter
\@ifundefined{KOMAClassName}{% if non-KOMA class
  \IfFileExists{parskip.sty}{%
    \usepackage{parskip}
  }{% else
    \setlength{\parindent}{0pt}
    \setlength{\parskip}{6pt plus 2pt minus 1pt}}
}{% if KOMA class
  \KOMAoptions{parskip=half}}
\makeatother
\usepackage{longtable,booktabs,array}
\usepackage{caption}
\captionsetup[table]{skip=6pt}
\newcounter{none} % for unnumbered tables
\usepackage{calc} % for calculating minipage widths
\usepackage{caption}
% Make caption package work with longtable
\makeatletter
\def\fnum@table{\tablename~\thetable}
\makeatother
\setlength{\emergencystretch}{3em} % prevent overfull lines
\providecommand{\tightlist}{%
  \setlength{\itemsep}{0pt}\setlength{\parskip}{0pt}}
% Auto-generated by infrastructure/rendering/_pdf_latex_helpers.py::extract_math_font_preamble.
% Minimal header passed to Pandoc via -H so Beamer slides pick up the same math font
% as the combined-PDF path. Adjust the manuscript's preamble.md to change the font globally.
\usepackage{unicode-math}
\setmathfont{latinmodern-math.otf}

% Manuscript macro/package fallbacks (newcommand -> providecommand).
% Core mathematics
\usepackage{amsmath}
\usepackage{amssymb}
% Document layout
% Tables
\usepackage{booktabs}
% Caption styling (booktabs already loaded above). Small caption font with a
% bold label ("Figure 1", "Table 2") gives the math/figure-heavy layout a
% consistent, journal-like caption rhythm without fighting pandoc-crossref —
% pandoc emits standard \caption{}, and caption/captionsetup only restyle it.
% ── Theorem environments ─────────────────────────────────────────────────
% amsthm is loaded above. The FedGVI<->Friston bridge is stated as a small
% number of theorems/definitions/lemmas/propositions/corollaries; sharing one
% counter (definition/lemma/proposition/corollary all step the `theorem`
% counter) gives a single monotone numbering 1, 2, 3, ... across all kinds, so
% authors can reference them as Theorem 1, Definition 2, Corollary 3 without a
% per-kind counter drifting out of order. Plain style for theorem-like results,
% definition style (upright body) for definitions.
% (algorithm2e intentionally not loaded: this manuscript states results as
% numbered amsthm Theorems/Definitions, not algorithm2e pseudocode.)
% Code listings
\usepackage{listings}
% Typography and formatting
% (siunitx intentionally not loaded: no \SI/\num units appear in this manuscript.)
% Cross-references and citations
% (cleveref and titlesec intentionally not loaded: unavailable in this TeX tree
% and not required. Cross-references resolve through pandoc-crossref's own
% \ref-style names — no \cref appears — and section heading styling falls back
% to the pandoc/LaTeX defaults.)
% ── TeX Live Basic-compatible mono font for code listings ────────────
% Latin Modern Mono is bundled with TeX Live Basic and keeps the manuscript
% renderable on minimal TeX installations. Mathematical Unicode coverage is
% provided separately by Latin Modern Math below, so the code font does not
% need a private JuliaMono installation.
%
% Use the TeX font filename rather than the system-family name: the minimal
% TeX Live fontconfig database may not register the OpenType family even though
% kpsewhich can resolve the bundled files.
% Math font for unicode-math: Latin Modern Math (TeX Live) has full BMP
% coverage including U+2223 (\mid), U+226A/226B (\ll/\gg), and the Greek/
% blackboard letters used in equations. Without an explicit \setmathfont,
% unicode-math falls back to lmroman text font which lacks several glyphs
% and emits "Missing character" warnings on every \mid in math mode.

% Slide layout defaults for warning-clean scientific decks.
\usepackage{etoolbox}
\IfFileExists{xurl.sty}{\usepackage{xurl}}{}
\IfFileExists{seqsplit.sty}{\usepackage{seqsplit}}{\newcommand{\seqsplit}[1]{#1}}
\protected\def\breakseq#1{\seqsplit{#1}}
\protected\def\breaktt#1{\begingroup\ttfamily\seqsplit{#1}\endgroup}
\setlength{\emergencystretch}{6em}
\tolerance=5000
\hbadness=10000
\hfuzz=1pt
\setlength{\tabcolsep}{2pt}
\AtBeginEnvironment{longtable}{\tiny\renewcommand{\arraystretch}{0.86}\setlength{\tabcolsep}{1pt}}
\AtBeginEnvironment{tabular}{\tiny\renewcommand{\arraystretch}{0.86}\setlength{\tabcolsep}{1pt}}
\AtBeginEnvironment{equation}{\tiny}
\AtBeginEnvironment{equation*}{\tiny}
\AtBeginEnvironment{align}{\tiny}
\AtBeginEnvironment{align*}{\tiny}
\AtBeginEnvironment{itemize}{\footnotesize}
\AtBeginEnvironment{enumerate}{\footnotesize}
\AtBeginEnvironment{description}{\footnotesize}
\setbeamerfont{caption}{size=\tiny}
\setbeamerfont{caption name}{size=\tiny}
\setbeamerfont{normal text}{size=\small}
\setbeamerfont{frametitle}{size=\small}
\setbeamerfont{section title}{size=\footnotesize}
\setbeamerfont{subsection title}{size=\footnotesize}
\setbeamertemplate{section page}{%
  \centering
  \begin{beamercolorbox}[sep=12pt,center,wd=\paperwidth]{section title}
    \parbox{0.86\paperwidth}{\centering\usebeamerfont{section title}\insertsection\par}
  \end{beamercolorbox}
}
\setbeamertemplate{subsection page}{%
  \centering
  \begin{beamercolorbox}[sep=8pt,center,wd=\paperwidth]{subsection title}
    \parbox{0.86\paperwidth}{\centering\usebeamerfont{subsection title}\insertsubsection\par}
  \end{beamercolorbox}
}
\setlength{\abovecaptionskip}{2pt}
\setlength{\belowcaptionskip}{0pt}

% Natbib and cross-reference fallbacks — slides don't load natbib
% or cleveref, but combined-PDF manuscript prose may emit these
% commands. The fallback renders citations as a bracketed key list
% and cross-references as detokenized labels so slides don't fail on
% undefined control sequences or raw underscores. \providecommand is
% a no-op if packages load later.
\providecommand{\citep}[1]{[#1]}
\providecommand{\citet}[1]{#1}
\providecommand{\citealp}[1]{#1}
\providecommand{\citeauthor}[1]{#1}
\providecommand{\citeyear}[1]{#1}
\providecommand{\cref}[1]{\texttt{\detokenize{#1}}}
\providecommand{\Cref}[1]{\texttt{\detokenize{#1}}}
\providecommand{\crefrange}[2]{\texttt{\detokenize{#1}}--\texttt{\detokenize{#2}}}
\providecommand{\Crefrange}[2]{\texttt{\detokenize{#1}}--\texttt{\detokenize{#2}}}
\providecommand{\paragraph}[1]{\textbf{#1}\ }

% Opt-in accessible presentation profile.
\setbeamerfont{normal text}{size*={20pt}{24pt}}
\setbeamerfont{frametitle}{size*={28pt}{32pt}}
\setbeamerfont{section title}{size*={28pt}{32pt}}
\setbeamerfont{subsection title}{size*={28pt}{32pt}}
\setbeamerfont{caption}{size*={16pt}{19pt}}
\setbeamerfont{caption name}{size*={16pt}{19pt}}
\setbeamerfont{itemize/enumerate body}{size*={20pt}{24pt}}
\setbeamerfont{itemize/enumerate subbody}{size*={20pt}{24pt}}
\setbeamerfont{itemize/enumerate subsubbody}{size*={20pt}{24pt}}
\setbeamerfont{itemize item}{size*={20pt}{24pt}}
\setbeamerfont{itemize subitem}{size*={20pt}{24pt}}
\setbeamerfont{itemize subsubitem}{size*={20pt}{24pt}}
\setbeamerfont{description body}{size*={20pt}{24pt}}
\setbeamerfont{description item}{size*={20pt}{24pt}}
\setbeamerfont{quote}{size*={20pt}{24pt}}
\setbeamerfont{quotation}{size*={20pt}{24pt}}
\AtBeginEnvironment{longtable}{\fontsize{20pt}{24pt}\selectfont}
\AtBeginEnvironment{tabular}{\fontsize{20pt}{24pt}\selectfont}
\AtBeginEnvironment{equation}{\fontsize{20pt}{24pt}\selectfont}
\AtBeginEnvironment{equation*}{\fontsize{20pt}{24pt}\selectfont}
\AtBeginEnvironment{align}{\fontsize{20pt}{24pt}\selectfont}
\AtBeginEnvironment{align*}{\fontsize{20pt}{24pt}\selectfont}
\AtBeginEnvironment{itemize}{\fontsize{20pt}{24pt}\selectfont}
\AtBeginEnvironment{enumerate}{\fontsize{20pt}{24pt}\selectfont}
\AtBeginEnvironment{description}{\fontsize{20pt}{24pt}\selectfont}
\AtBeginDocument{\usebeamerfont{normal text}}
\newlength{\AFSlideTableWidth}
% The fixed footer reservation is calibrated with the semantic seven-line regular-body budget.
\setbeamertemplate{footline}{%
  \leavevmode\vbox{%
    \hbox{%
      \begin{beamercolorbox}[wd=\paperwidth,ht=16pt,dp=3pt,center]{author in head/foot}%
        \fontsize{16pt}{19pt}\selectfont Untagged PDF derivative \textbar\ \href{../web/index.html}{HTML reader}%
      \end{beamercolorbox}%
    }%
    \vskip6pt%
  }%
}
% Accessible footer reserved height: 25pt.

\newtheorem{proposition}{Proposition}
\newtheorem{hypothesis}{Hypothesis}
\newtheorem{remark}[theorem]{Remark}
\newtheorem{axiom}{Axiom}
\newtheorem{property}{Property}
\usepackage{bookmark}
\IfFileExists{xurl.sty}{\usepackage{xurl}}{} % add URL line breaks if available
\urlstyle{same}
% fallback for those not using the hyperref driver hyperxmp:
\makeatletter
\@ifundefined{xmpquote}{\newcommand{\xmpquote}[1]{#1}}{}
\makeatother
\hypersetup{
  hidelinks,
  pdfcreator={LaTeX via pandoc}}

\author{\texorpdfstring{}{}}
\date{}

\begin{document}

\begin{frame}{Contamination models: declared failure modes for belief
fusion}
\protect\phantomsection\label{sec:methods-contamination}
Robust belief fusion only earns its keep when some agents are wrong.
\end{frame}

\begin{frame}{Contamination models (part 2)}
The active-inference community has built ensembles that coordinate by
sharing beliefs and observations (Friston et al. 2024; Heins et al.
2024; Albarracin et al. 2022; Kaufmann et al. 2021),
\end{frame}

\begin{frame}{Contamination models (part 3)}
but it has assumed those beliefs are trustworthy in the cited modeled
protocols: fusion is treated as exact-Bayes pooling of well-calibrated
reports.
\end{frame}

\begin{frame}{Contamination models (part 4)}
The robust-Bayes and federated-learning literatures (McMahan et al.
2017; Ashman et al. 2022; Mildner et al. 2025) have, in turn, studied
robustness to corrupted clients under their declared settings, but
outside the generative-model-bearing POMDP setting.
\end{frame}

\begin{frame}{Contamination models (part 5)}
This section defines the corruption process that lets us test fusion
robustness inside the active-inference colony --- the experimental
complement of the robust aggregation rule of
sec.~5.8.
\end{frame}

\begin{frame}[fragile]{Corruption process for adversarial belief
broadcasts}
\protect\phantomsection\label{sec:methods-corruption}
In the sentinel world a healthy sentinel reports a well-calibrated
categorical over the creature's location; a contaminated one reports
something corrupted. \texttt{contamination.contaminate} manufactures the
corrupted reports.
\end{frame}

\begin{frame}{Corruption process for\ldots{} (part 2)}
Every corruption is a convex mixture of the agent's belief \(b\) with a
corruption target \(t\), governed by a single rate \(r \in [0, 1]\),
\end{frame}

\begin{frame}{Corruption process for\ldots{} (part 3)}
\begin{equation}\tag{21}\protect\phantomsection\label{eq:contamination-mix}{
\tilde b \;=\; (1 - r)\, b \;+\; r\, t,
}\end{equation}
\end{frame}

\begin{frame}{Corruption process for\ldots{} (part 4)}
so the experiments sweep exactly one knob. The convex form of
eq.~21 gives a clean limit and is the anchor of
the suite (ISC-26): at \(r = 0\) every corruption kind returns the input
belief unchanged, so contamination is a strict, continuous departure
from the uncorrupted Friston belief-share ---
\end{frame}

\begin{frame}{Corruption process for\ldots{} (part 5)}
never a discontinuity.
\end{frame}

\begin{frame}[fragile]{Corruption process for\ldots{} (part 6)}
This section defines the three core corruption targets \(t\), each
capturing a distinct failure of a federated agent. Geometrically the
three are three landmarks of the probability simplex --- a wrong vertex
(\texttt{confident\_wrong}), the flat centroid (\texttt{uniform}), and a
random interior point (\texttt{label\_noise}) ---
\end{frame}

\begin{frame}{Corruption process for\ldots{} (part 7)}
so the mixture of eq.~21 drags an honest belief
toward a qualitatively different destination in each case.
\end{frame}

\begin{frame}[fragile]{Corruption process for\ldots{} (part 8)}
Two further mechanisms (\texttt{byzantine} and \texttt{drift}) extend
the same convex-mix contract and are introduced in the extended-methods
supplement (sec.~13.2).
\end{frame}

\begin{frame}[fragile]{Corruption process for\ldots{} (part 9)}
\textbf{\texttt{confident\_wrong} --- the adversarial sentinel.} This is
the lookout that points to one wrong cell and insists on it with total
certainty. The target is a one-hot spike on a wrong state,
\(t = \mathrm{onehot}(s_{\text{wrong}})\), so \(\tilde b\) is mixed
toward a confident, mistaken delta.
\end{frame}

\begin{frame}{Corruption process for\ldots{} (part 10)}
Callers choose \(s_{\text{wrong}}\) explicitly; the verdict sweep of
sec.~5.22 fixes it once per colony as
the state diametrically opposite the true state on the location grid,
held constant across the entire rate sweep, rather than deriving it from
the agent's current belief.
\end{frame}

\begin{frame}{Corruption process for\ldots{} (part 11)}
At \(r = 1\) this is a pure delta on the wrong cell. This is the
saboteur that is \emph{sure} and \emph{mistaken}: exactly the agent that
robust aggregation must reject.
\end{frame}

\begin{frame}[fragile]{Corruption process for\ldots{} (part 12)}
\textbf{\texttt{label\_noise} --- the miscalibrated sentinel.} This is
the lookout with a scrambled sensor: it is not lying toward any
particular cell, only diluting every honest report with the same fixed
sprinkle of noise. The target is a fixed noisy categorical drawn once
from a \(\mathrm{Dirichlet}(1)\) (a random but valid pmf),
\end{frame}

\begin{frame}{Corruption process for\ldots{} (part 13)}
modeling a sentinel whose report is partly random rather than
adversarial.
\end{frame}

\begin{frame}{Corruption process for\ldots{} (part 14)}
Because the noisy target is drawn once and then held fixed across the
rate sweep, the corruption has no direction to exploit and no single
cell to veto --- the robust pool meets diffuse degradation, not a
targeted attack.
\end{frame}

\begin{frame}[fragile]{Corruption process for\ldots{} (part 15)}
\textbf{\texttt{uniform} --- the apathetic sentinel.} This is the
lookout that shrugs: it has lost track of the creature and calls every
cell equally likely. The target is the maximum-entropy uniform pmf
\(t = (1/n_s)\mathbf{1}\), modeling a saturated sentinel that has lost
all information. At \(r = 1\) it reports uniform,
\end{frame}

\begin{frame}{Corruption process for\ldots{} (part 16)}
contributing no evidence to the pool rather than actively pulling it
toward a wrong cell.
\end{frame}

\begin{frame}[fragile]{Corruption process for\ldots{} (part 17)}
All three share the contract of eq.~21 and
require an explicit seeded generator --- \texttt{label\_noise} uses it
to draw the noisy target --- so every contaminated report is
reproducible. The grid of rates the sweep uses,
\(\{0, 0.225, 0.45, 0.675, 0.9\}\), deliberately stops below the
pure-veto limit \(r = 1\),
\end{frame}

\begin{frame}{Corruption process for\ldots{} (part 18)}
where a fully-confident wrong delta forces \emph{every} pooling rule's
accuracy to zero and the robust-versus-naive contrast vanishes.
\end{frame}

\begin{frame}{How contamination meets the three robustness axes}
\protect\phantomsection\label{sec:methods-contamination-axes}
The authoritative guarantee taxonomy is sec.~3.3;
this section only records where the declared corruption enters each
route.
\end{frame}

\begin{frame}{Client bounded-loss route}
\protect\phantomsection\label{sec:methods-contamination-client}
A corrupted local observation enters the per-agent generalized-Bayes
update of sec.~5.7. FedGVI's bounded-influence
result applies only under the source theorem's matching loss,
divergence, model, and regularity assumptions (Mildner et al. 2025).
\end{frame}

\begin{frame}[fragile]{Client bounded-loss route (part 2)}
The logistic-regression comparison in fig.~18 is a
narrower exploratory proxy: it compares RCCE and NLL point-estimate
gradients under two L2 coefficients. Its legacy \texttt{AR} argument
selects the larger coefficient;
\end{frame}

\begin{frame}{Client bounded-loss route (part 3)}
it does not evaluate Alpha-Rényi divergence or inherit the source
theorem from the observed curve.
\end{frame}

\begin{frame}[fragile]{Heuristic server route}
\protect\phantomsection\label{sec:methods-contamination-heuristic-server}
A corrupted posterior broadcast enters \texttt{robust\_aggregate} at
fusion. Divergence reweighting can assign a smaller effective weight to
a report far from the current consensus, as the finite diagnostics
figs.~13, 15 illustrate.
\end{frame}

\begin{frame}{Heuristic server route (part 2)}
The rule's positive formal property is the exact \(c=0\) recovery limit
eq.~10, not a client-loss bound, Byzantine
guarantee, or universal robustness theorem.
\end{frame}

\begin{frame}[fragile]{Objective-backed variational server route}
\protect\phantomsection\label{sec:methods-contamination-variational-server}
The same corrupted posterior can instead enter
\texttt{variational\_aggregate}, whose declared free energy,
block-descent property, and effective-weight bound are given in
sec.~5.8.2.
\end{frame}

\begin{frame}[fragile]{Objective-backed variational\ldots{} (part 2)}
Those properties belong to that forward-oriented objective and do not
certify \texttt{robust\_aggregate} or the source client's bounded-loss
theorem.
\end{frame}

\begin{frame}{Objective-backed variational\ldots{} (part 3)}
Contamination is therefore a common stressor, not a license to transfer
a guarantee or empirical result between client, heuristic-server, and
variational-server lanes.
\end{frame}

\end{document}
